% Direct-authority diplomatic French for printed p.178.
\SourcePage{178}

\LayerTarget{tohoku:criterion:cech-H2-example-reduction:p2}
\noindent\textit{isomorphisme.}

\LayerTarget{tohoku:example:cech-H2-counterexample-construction}
Soit $X$ un espace irréductible (cf. fin de 3.3), $Y_1$ et $Y_2$ deux parties
fermées irréductibles de $X$ se rencontrant en exactement deux points $x_1$ et
$x_2$ (par exemple deux cercles sécants dans le plan muni de la topologie de
Zariski), $Y$ leur réunion. Par abus de notation, on désigne par $\mathbf Z$ le
faisceau constant des entiers sur $X$, et on considère, avec les notations de
3.5, la suite exacte de faisceaux

\LayerTarget{tohoku:eq:closed-complement-constant-sheaf}
\[
0\longrightarrow\mathbf Z_{C_Y}\longrightarrow\mathbf Z
\longrightarrow\mathbf Z_Y\longrightarrow 0.
\]

\LayerTarget{tohoku:claim:cech-H2-counterexample-sheaf}
Je dis que le faisceau $F=\mathbf Z_{C_Y}$ satisfait aux conditions voulues.
Tout d'abord, comme $H^p(X,\mathbf Z)=0$ pour $p>0$ d'après la fin de 3.3, on a
$H^2(X,\mathbf Z_{C_Y})=H^1(X,\mathbf Z_Y)=H^1(Y,\mathbf Z)$, montrons que ce
groupe n'est pas nul, et de façon précise est isomorphe à $\mathbf Z$. En
effet, on a un monomorphisme naturel du faisceau constant $\mathbf Z$ sur $Y$
dans la somme directe des faisceaux $\mathbf Z_{Y_1}$ et $\mathbf Z_{Y_2}$,
d'où une suite exacte de faisceaux sur $Y$ :

\LayerTarget{tohoku:eq:union-components-constant-sheaves}
\[
0\longrightarrow\mathbf Z_Y
\longrightarrow(\mathbf Z_{Y_1}+\mathbf Z_{Y_2})
\longrightarrow\mathbf Z_{Y_1\cap Y_2}\longrightarrow 0.
\]

\LayerTarget{tohoku:calculation:H1-union-cokernel}
Or $Y_i$ étant irréductible, on a encore
$H^p(Y,\mathbf Z_{Y_i})=H^p(Y_i,\mathbf Z)=0$ pour $p>0$, d'où
\[
H^1(Y,\mathbf Z)
=\Gamma(\mathbf Z_{Y_1\cap Y_2})/
\operatorname{Im}\bigl(\Gamma(\mathbf Z_{Y_2})\bigr),
\]
c'est le conoyau du homomorphisme de groupes
$\mathbf Z^2\longrightarrow\mathbf Z^2$ donné par
$(n_1,n_2)\longmapsto(n_1-n_2,n_1-n_2)$, c'est-à-dire un groupe isomorphe à
$\mathbf Z$, d'où

\LayerTarget{tohoku:eq:counterexample-H2-is-Z}
\[
H^2(X,\mathbf Z_{C_Y})=H^1(X,\mathbf Z)=\mathbf Z.
\]

\LayerTarget{tohoku:goal:counterexample-cech-H1-is-Z}
Il reste à prouver que
$\check H^1(X,\check H^1(\mathbf Z_{C_Y}))$ est isomorphe à $\mathbf Z$ (car
compte tenu de la relation précédente, l'épimorphisme
$\check H^2(X,F)\longrightarrow H^1(X,H^1(F))$ sera nécessairement un
isomorphisme). Calculons $H^1(\mathbf Z_{C_Y})$ ; pour tout ouvert $V$,
$H^1(V,\mathbf Z_{C_Y})$ peut se calculer grâce à la suite exacte suivante de
faisceaux sur $V$ :

\LayerTarget{tohoku:eq:restricted-complement-constant-sheaf}
\[
0\longrightarrow\mathbf Z_{C_Y}\longrightarrow\mathbf Z
\longrightarrow\mathbf Z_{Y'}\longrightarrow 0,
\]
où on a posé $Y'=Y\cap V$. Comme $V$ est aussi irréductible, on en conclut

\LayerTarget{tohoku:eq:local-H1-reduced-H0}
\[
H^1(V,\mathbf Z_{C_Y})
=H^0(Y',\mathbf Z)/\operatorname{Im}H^0(V,\mathbf Z)
=\widetilde H^0(Y',\mathbf Z),
\]
où le dernier groupe désigne le groupe de cohomologie entière réduit de
dimension 0, i.e. ici le groupe abélien libre engendré par les composantes
connexes de $Y'$, modulo le sous-groupe diagonal.

\LayerTarget{tohoku:criterion:local-H1-two-components}
Ici $Y'=Y\cap V$ est un sous-espace ouvert de $Y=Y_1\cup Y_2$, et a donc 0, 1
ou 2 composantes connexes, le dernier cas se présentant exactement si $V$
rencontre à la fois $Y_1$ et $Y_2$ sans rencontrer leur intersection ; donc on
a $H^1(V,\mathbf Z_{C_Y})=0$, sauf dans ce dernier cas.

\LayerTarget{tohoku:construction:special-point-indexed-coverings}
Pour calculer
\[
H^1(X,H^1(F))=\varinjlim H^1(\mathcal U,H^1(F)),
\]
on peut se limiter aux recouvrements $\mathcal U=(U_x)_{x\in X}$ tels que
chaque $U_x$ rencontre au plus un seul des deux fermés $Y_1$ et $Y_2$, sauf si
$x$ est l'un des deux points $x_1,x_2$ de $Y_1\cap Y_2$ auquel cas on suppose
qu'il ne contient pas l'autre.

\LayerTarget{tohoku:eq:special-cover-C0-vanishing}
Pour un tel $\mathcal U$, on voit aussitôt
$C^0(\mathcal U,H^1(F))=0$, donc $H^1(\mathcal U,H^1(F))$ s'identifie au groupe
$Z^1(\mathcal U,H^1(F))$ des 1-cocycles $\mathcal U$ à coefficients dans
$H^1(F)$, soient $(f_{x,y})_{x,y\in X}$.

\LayerTarget{tohoku:eq:special-cover-intersection-sections}
Mais on a
\[
\Gamma(U_x\cap U_y,H^1(F))=0
\]
sauf si $x=x_1$, $y=x_2$ ou $x=x_2$, $y=x_1$. D'où

\LayerTarget{tohoku:eq:special-cover-C1-is-Z2}
\[
C^1(\mathcal U,H^1(F))\simeq\mathbf Z^2,
\]
et on voit aussitôt, que les cocyles s'identifient aux couples $(n,-n)$
d'entiers opposés ($n=f_{x_1,x_2}$).

\LayerTarget{tohoku:eq:special-cover-H1-is-Z}
On a donc $H^1(\mathcal U,H^1(F))=\mathbf Z$, d'où aussitôt à

\TohokuContinuation{179}
